\section{Illustrations of Galois Theory}

\exercise{54.1}{Irreducibility of the eighth cyclotomic polynomial}
Make the invertible substitution $x=y+1$. Then
\[(y+1)^4+1=y^4+4y^3+6y^2+4y+2.\]
Eisenstein's criterion at 2 applies: every coefficient except the leading one is divisible by 2, and the constant 2 is not divisible by 4. Thus the substituted polynomial is irreducible over $\Q$. A factorization of $x^4+1$ would yield a factorization after substitution, so $x^4+1$ is irreducible too.

\exercise{54.2}{Checking the entire octic fixed-field diagram}
Put $\alpha=\sqrt[4]2$ and $K=\Q(\alpha,i)$. We have $[K:\Q]=8$, since $x^4-2$ is Eisenstein and $\Q(\alpha)$ is real. To avoid relying on a picture alone, write $r(\alpha)=i\alpha,r(i)=i$ and $s(\alpha)=\alpha,s(i)=-i$. Then $r^4=s^2=1$ and $srs=r^{-1}$. In the source notation $r=\rho_1$, $r^2=\rho_2$, $s=\mu_1$, $r^2s=\mu_2$, $rs=\delta_1$, and $r^3s=\delta_2$.
The proposed fields can all be verified by their fixed generators and degrees:
\[
\begin{array}{c|c|c}
\text{subgroup}&\text{proposed fixed field}&\text{degree over }\Q\\\hline
\langle r^2,s\rangle=H_1&\Q(\alpha^2)=\Q(\sqrt2)&2\\
\langle r\rangle=H_2&\Q(i)&2\\
\langle r^2,rs\rangle=H_3&\Q(i\alpha^2)=\Q(\sqrt{-2})&2\\
\langle s\rangle=H_4&\Q(\alpha)&4\\
\langle r^2s\rangle=H_5&\Q(i\alpha)&4\\
\langle r^2\rangle=H_6&\Q(\alpha^2,i)&4\\
\langle rs\rangle=H_7&\Q((1+i)\alpha)&4\\
\langle r^3s\rangle=H_8&\Q((1-i)\alpha)&4
\end{array}
\]
For each row, substitution of $r,s$ checks that the indicated generator is fixed by the subgroup. For $H_1,H_2,H_3$, the generator is a nonrational quadratic radical, giving degree 2. For $H_4,H_5$, the fourth power is 2 and $x^4-2$ is irreducible. For $H_6$, the real quadratic field gains the nonreal $i$, giving degree 4. For $H_7,H_8$, if $\gamma=(1\pm i)\alpha$, then $\gamma^2=\pm2i\sqrt2$ and $\gamma^4=-8$. Its square generates $\Q(\sqrt{-2})$, and $\pm2i\sqrt2$ is not a square there: writing $(a+b\sqrt{-2})^2=\pm2\sqrt{-2}$ would force $a^2-2b^2=0$ and $ab=\pm1$, impossible with rational $a,b$. Thus these fields have degree 4 too. In every row this degree equals the subgroup index, proving equality with the full fixed field. The trivial subgroup fixes $K$; the full group fixes $\Q$. This verifies all ten fields and, by reversing subgroup containment, all inclusions in the diagram.

\exercise{54.3}{Primitive generators and all their minimal polynomials}
Use the notation of Exercise 54.2. A complete list, including top and bottom, is
\[
\begin{array}{c|c|c}
\text{field}&\text{primitive element}&\text{minimal polynomial}\\\hline
\Q&0&x\\
\Q(\sqrt2)&\sqrt2&x^2-2\\
\Q(i)&i&x^2+1\\
\Q(\sqrt{-2})&i\sqrt2&x^2+2\\
\Q(\alpha)&\alpha&x^4-2\\
\Q(i\alpha)&i\alpha&x^4-2\\
\Q(\sqrt2,i)&\sqrt2+i&x^4-2x^2+9\\
\Q((1+i)\alpha)&(1+i)\alpha&x^4+8\\
\Q((1-i)\alpha)&(1-i)\alpha&x^4+8\\
K&\alpha+i&x^8+4x^6+2x^4+28x^2+1
\end{array}
\]
The degrees and irreducibility in the first nine rows follow from Exercise 54.2 and the four distinct sign conjugates of $\sqrt2+i$. For the last row, the eight images of $\alpha+i$ are $i^j\alpha+\varepsilon i$ for $0\le j<4$ and $\varepsilon=\pm1$. They are all distinct: an equality with different $j$ would express $\alpha$ as an element of $\Q(i)$, impossible because its rational degree is 4; with equal $j$ the signs must agree. Thus this element has degree 8 and generates $K$. Multiplying
\[((x-i)^4-2)((x+i)^4-2)\]
gives the displayed monic polynomial, proving it is its minimal polynomial. The identical quartics in several rows do not imply identical subfields: an abstract isomorphism can send one primitive root to another.

\exercise{54.4}{The fifth-root-of-unity correspondence}
Let $\zeta$ be a primitive fifth root.
\begin{enumerate}[label=(\alph*)]
\item Its powers $1,\zeta,\zeta^2,\zeta^3,\zeta^4$ are all roots of $x^5-1$, so $K=\Q(\zeta)$ is exactly its splitting field.
\item An automorphism preserves order and the equation $\zeta^5=1$, so sends $\zeta$ to $\zeta^a$ for $a=1,2,3,4$.
\item $\Phi_5$ is irreducible of degree 4. Thus every listed image gives an automorphism, $\sigma_a(\zeta)=\zeta^a$. Composition multiplies exponents modulo 5. Since 2 has order 4 modulo 5, the group is $\langle\sigma_2\rangle\cong\Z_4$.
\item Its only proper nontrivial subgroup is $\langle\sigma_4\rangle$, generated by complex conjugation. Its fixed field is $M=\Q(t)$ where $t=\zeta+\zeta^{-1}$. Dividing $1+\zeta+\cdots+\zeta^4=0$ by $\zeta^2$ gives $t^2+t-1=0$, so $M=\Q(\sqrt5)$ has degree 2 and is the full fixed field. The two diagrams are the chains
\[
\begin{array}{c@{\qquad\qquad}c}
K&\langle\sigma_2\rangle\\
\vert&\vert\\
M&\langle\sigma_4\rangle\\
\vert&\vert\\
\Q&\{1\}
\end{array}
\]
with correspondence $K\leftrightarrow\{1\}$, $M\leftrightarrow\langle\sigma_4\rangle$, and $\Q\leftrightarrow\langle\sigma_2\rangle$. Each step has index or field degree 2.
\end{enumerate}

\exercise{54.5}{A fifth radical over a cyclotomic base}
Put $F=\Q(\zeta_5)$ and $r=\sqrt[5]2$. Since $[F:\Q]=4$ and $[\Q(r):\Q]=5$, the tower law applied to $F(r)$ shows $5\mid4[F(r):F]$. But adjoining a root of $x^5-2$ gives degree at most 5, so $[F(r):F]=5$. The presence of $\zeta_5$ makes all roots $r\zeta_5^j$ lie in $F(r)$, so this is the splitting field. The map $r\mapsto r\zeta_5$, fixing $F$, has order 5 and generates all five automorphisms. Thus the group is $\Z_5$.

\exercise{54.6}{The seventh-root-of-unity correspondence}
Let $\zeta$ have order 7 and $K=\Q(\zeta)$.
\begin{enumerate}[label=(\alph*)]
\item Its seven powers are precisely the roots of $x^7-1$, so it generates the splitting field.
\item Every automorphism sends $\zeta$ to one of $\zeta^a$, $1\le a\le6$.
\item Irreducibility of $\Phi_7$ gives degree 6, so all six exponent maps occur. Their group is $(\Z/7\Z)^\times$, cyclic of order 6 generated by 3. Write $g=\sigma_3$.
\item The two proper nontrivial subgroups are $\langle g^2\rangle=\{\sigma_1,\sigma_2,\sigma_4\}$ of order 3 and $\langle g^3\rangle=\{\sigma_1,\sigma_6\}$ of order 2. Their fixed fields are respectively $A=\Q(\sqrt{-7})$ and $B=\Q(\zeta+\zeta^{-1})$. For the first, put $\eta=\zeta+\zeta^2+\zeta^4$. The other three powers sum to $\bar\eta$, and direct multiplication gives $\eta+\bar\eta=-1$, $\eta\bar\eta=2$. Thus $\eta$ satisfies $x^2+x+2$, has degree 2, and generates $A$. The subgroup fixes this sum, so it fixes exactly $A$. For $B$, complex conjugation fixes $t=\zeta+\zeta^{-1}$ and $\zeta$ satisfies $X^2-tX+1$; since $B$ is real and $\zeta$ nonreal, $[K:B]=2$ and $[B:\Q]=3$. For a check, $t^3+t^2-2t-1=0$ follows by dividing $\Phi_7(\zeta)=0$ by $\zeta^3$.
The field and group diagrams are diamonds:
\begin{center}\setlength{\unitlength}{1pt}
\begin{picture}(360,130)
\put(80,120){\makebox(0,0){$K$}}\put(20,65){\makebox(0,0){$A$}}\put(140,65){\makebox(0,0){$B$}}\put(80,10){\makebox(0,0){$\Q$}}
\Line(74,110)(26,75)\Line(86,110)(134,75)\Line(26,55)(74,20)\Line(134,55)(86,20)
\put(280,120){\makebox(0,0){$\langle g\rangle$}}\put(220,65){\makebox(0,0){$\langle g^2\rangle$}}\put(340,65){\makebox(0,0){$\langle g^3\rangle$}}\put(280,10){\makebox(0,0){$\{1\}$}}
\Line(274,110)(226,75)\Line(286,110)(334,75)\Line(226,55)(274,20)\Line(334,55)(286,20)
\end{picture}\end{center}
The correspondence reverses order: $A\leftrightarrow\langle g^2\rangle$, $B\leftrightarrow\langle g^3\rangle$, with top and bottom paired oppositely. There are no other intermediate fields because a cyclic group of order 6 has no other subgroups.
\end{enumerate}

\exercise{54.7}{Eighth roots without recomputing a table}
The splitting field is $\Q(\zeta_8)=\Q(\sqrt2,i)$. Its group is $(\Z/8\Z)^\times=\{1,3,5,7\}$ under multiplication modulo 8. Each of its three nonidentity elements squares to 1, so it is the Klein four-group $\Z_2\times\Z_2$.

\exercise{54.8}{A complete quartic splitting-field correspondence}
\step{Find the field and its degree.} Put $a=\sqrt{2+\sqrt5}$, $b=i/a$, and $K=\Q(a,i)$. Then $b^2=2-\sqrt5$ and $ab=i$, so the roots are $a,-a,b,-b$ and $K$ is their splitting field. The norm of $2+\sqrt5$ from $\Q(\sqrt5)$ to $\Q$ is $-1$, so it is not a square there. Thus $[\Q(a):\Q]=4$; this field is real and gains degree 2 on adjoining $i$, so $[K:\Q]=8$.
\step{Describe the group.} Define
\[r(a)=b,\quad r(i)=-i,\qquad s(a)=a,\quad s(i)=-i.\]
To justify existence, conjugacy gives an embedding $\Q(a)\to\Q(b)$ sending $a$ to $b$. The element $i$ is outside $\Q(b)$: otherwise $a=i/b$ would also belong to that degree-four field, contradicting $[K:\Q]=8$. Thus $x^2+1$ is still irreducible over $\Q(b)$, and the embedding extends with either choice $i\mapsto\pm i$. The map $s$ is ordinary complex conjugation. This constructs the asserted automorphisms, rather than merely proposing images. They satisfy $r^4=s^2=1$, $srs=r^{-1}$, and yield eight distinct maps $1,r,r^2,r^3,s,rs,r^2s,r^3s$. Thus the group is the dihedral group of order 8.
\step{List every subgroup and its field.} Define the field labels by this table:
\[
\begin{array}{c|c|c}
\text{subgroup}&\text{fixed field}&[\text{field}:\Q]\\\hline
G&\Q&1\\
\langle r\rangle&B=\Q(i\sqrt5)&2\\
\langle r^2,s\rangle&A=\Q(\sqrt5)&2\\
\langle r^2,rs\rangle&C=\Q(i)&2\\
\langle r^2\rangle&E_0=\Q(\sqrt5,i)&4\\
\langle s\rangle&E_1=\Q(a)&4\\
\langle r^2s\rangle&E_2=\Q(ia)&4\\
\langle rs\rangle&E_3=\Q(a+b)&4\\
\langle r^3s\rangle&E_4=\Q(a-b)&4\\
\{1\}&K&8
\end{array}
\]
For verification, $r$ negates both $i$ and $\sqrt5=a^2-2$, hence fixes $i\sqrt5$. The subgroup $\langle r^2,s\rangle$ fixes $a^2$, and $\langle r^2,rs\rangle$ fixes $i$. The individual involutions fix their displayed generators: $r^2$ fixes $a^2,i$; $s$ fixes $a$; $r^2s$ fixes $ia$; $rs$ interchanges $a,b$; $r^3s$ interchanges $a,-b$.
The quadratic degrees are immediate. The field $E_1=\Q(a)$ has degree 4 as proved above. For $E_2$, the element $ia$ satisfies $X^4+4X^2-1=0$. Its images under the already constructed automorphisms include $ia,-ia,ib,-ib$, four distinct numbers: $\pm ia$ are nonzero imaginary and $\pm ib=\mp1/a$ are nonzero real. Its minimal polynomial therefore has at least four distinct roots and degree at most 4, so its degree is exactly 4. For $E_0$, $i$ doubles the real quadratic degree. Finally $(a\pm b)^2=4\pm2i$, which is not a square in $\Q(i)$: the norm of a putative square would be the rational square 20, impossible. Hence $E_3,E_4$ have degree 4, with minimal polynomial $x^4-8x^2+20$. All proposed fields therefore have the required index and equal the full fixed fields.
\step{Give all inclusions.} The following is the full field diagram. Reverse its edges and replace each field by its subgroup from the table to obtain the full subgroup diagram. In particular $E_0$ contains all three quadratic fields; $E_1,E_2$ contain only $A$, and $E_3,E_4$ contain only $C$.
\begin{center}\setlength{\unitlength}{1pt}
\begin{picture}(420,200)
\put(210,190){\makebox(0,0){$K$}}
\put(30,130){\makebox(0,0){$E_1$}}\put(115,130){\makebox(0,0){$E_2$}}\put(210,130){\makebox(0,0){$E_0$}}\put(305,130){\makebox(0,0){$E_3$}}\put(390,130){\makebox(0,0){$E_4$}}
\put(95,70){\makebox(0,0){$A$}}\put(210,70){\makebox(0,0){$B$}}\put(325,70){\makebox(0,0){$C$}}\put(210,10){\makebox(0,0){$\Q$}}
\Line(200,182)(35,140)\Line(205,180)(118,140)\Line(210,180)(210,140)\Line(215,180)(302,140)\Line(220,182)(385,140)
\Line(35,120)(90,80)\Line(112,120)(98,80)\Line(205,120)(100,80)\Line(210,120)(210,80)\Line(215,120)(320,80)\Line(308,120)(322,80)\Line(385,120)(330,80)
\Line(100,60)(200,18)\Line(210,60)(210,20)\Line(320,60)(220,18)
\end{picture}\end{center}
To make the subgroup diagram explicit as well, its covering relations are
\[
\begin{array}{c|l}
\text{upper subgroup}&\text{immediately lower subgroups}\\\hline
G&\langle r\rangle,\ \langle r^2,s\rangle,\ \langle r^2,rs\rangle\\
\langle r\rangle&\langle r^2\rangle\\
\langle r^2,s\rangle&\langle r^2\rangle,\langle s\rangle,\langle r^2s\rangle\\
\langle r^2,rs\rangle&\langle r^2\rangle,\langle rs\rangle,\langle r^3s\rangle\\
\text{each order-two subgroup}&\{1\}
\end{array}
\]
These relations and the field diagram exhibit the entire correspondence, not just its normal subgroups.

\exercise{54.9}{Two symmetric-function identities}
Write $s_1=y_1+y_2+y_3$, $s_2=y_1y_2+y_1y_3+y_2y_3$, $s_3=y_1y_2y_3$.
\begin{enumerate}[label=(\alph*)]
\item Expanding $s_1^2$ gives $y_1^2+y_2^2+y_3^2+2s_2$, so the answer is $s_1^2-2s_2$.
\item Put the six displayed ratios over the denominator $s_3$. Their numerator is $\sum_{i\ne j}y_i^2y_j$. On expanding $s_1s_2$, this sum appears together with $3y_1y_2y_3$. Thus the answer is $(s_1s_2-3s_3)/s_3$. The equality is an identity of rational functions, and numerical evaluation requires each original denominator nonzero.
\end{enumerate}

\exercise{54.10}{Polynomials of a sum and of squares}
Vieta's identities for $x^3-4x^2+6x-2$ give $s_1=4,s_2=6,s_3=2$ for its three roots.
\begin{enumerate}[label=(\alph*)]
\item Their sum is 4, so the monic polynomial with that sole zero is $x-4$.
\item The sum of the squared roots is $s_1^2-2s_2=4$. Their pairwise products sum to $s_2^2-2s_1s_3=36-16=20$, and their product is $s_3^2=4$. Hence the required polynomial is $x^3-4x^2+20x-4$. These squared roots are distinct: an opposite pair of original roots would force the third root to be 4, but the original polynomial does not vanish at 4.
\end{enumerate}

\exercise{54.11}{Every finite group occurs over a suitable field}
Let $H$ be finite. By Cayley's theorem embed it into $S_n$ for some $n$. Take independent indeterminates $y_1,\ldots,y_n$ over $\Q$, and put $L=\Q(y_1,\ldots,y_n)$ and $B=\Q(s_1,\ldots,s_n)$, using their elementary symmetric functions. The symmetric-function theorem gives $G(L/B)\cong S_n$. Regard $H$ as a subgroup and let $F=L_H$. The Galois correspondence gives $G(L/F)=H$, and $L/F$ is a finite normal extension. This proves existence over a suitable field $F$; it does not assert that every finite group occurs over $\Q$.

\exercise{54.12}{Discriminants and even permutations}
Let $\Delta=\prod_{i<j}(\alpha_i-\alpha_j)$, with roots repeated according to the multiplicities in $f$.
\begin{enumerate}[label=(\alph*)]
\item $\Delta=0$ exactly when two entries in this root list coincide, meaning $f$ has a repeated root. Since each irreducible factor is separable, repetitions occur exactly when some irreducible factor appears at least twice. This is the stated square-divisibility criterion.
\item An automorphism permutes the root list and sends $\Delta$ to $\pm\Delta$. Squaring removes the sign, so $\Delta^2$ is fixed by the full Galois group and belongs to $F$.
\item The equivalence needs two additional hypotheses: $f$ must have distinct roots and $\operatorname{char}F\ne2$. Under those hypotheses the root action embeds $G(K/F)$ into $S_n$, and
\[\sigma(\Delta)=\operatorname{sgn}(\sigma)\Delta.\]
This sign rule follows because interchanging two roots negates the alternating product and transpositions generate $S_n$. Since $\Delta\ne0$ and $1\ne-1$, it is fixed exactly by even permutations. Thus $G(K/F)\subseteq A_n$ if and only if every automorphism fixes $\Delta$, equivalently $\Delta\in F$.
As printed, the unrestricted assertion is false. In characteristic 2, $x^2+x+1$ over $\Z_2$ has splitting field of order 4 and group $S_2$, but $\Delta=\alpha-\alpha^2=1\in\Z_2$. With repeated roots, $\Delta=0\in F$ regardless of the group, and the set of distinct roots may have fewer than $n$ elements, so the stated identification with permutations of $n$ distinct roots is not available. The preceding corrected version is the usual discriminant test.
\end{enumerate}

\exercise{54.13}{Algebraic integers form a ring}
Let $\alpha,\beta$ satisfy monic integer polynomials of degrees $m,n$. Higher powers of either can be reduced using its monic relation with integer coefficients. Thus the additive group
\[M=\sum_{0\le j<m,\ 0\le k<n}\Z\alpha^j\beta^k\]
is finitely generated, contains 1, and is closed under multiplication by $\alpha$ and $\beta$, hence by $\alpha+\beta$, $\alpha\beta$, and $-\alpha$. It is torsion free as a subgroup of the additive group of $\C$, so it is a free abelian group of finite rank. For each $\gamma$ among these three elements, multiplication by $\gamma$ on $M$ is represented in a $\Z$-basis by an integer matrix $A$. Its characteristic polynomial is monic with integer coefficients, and Cayley--Hamilton gives $p(A)=0$. Applying that identity to $1\in M$ yields $p(\gamma)=0$ in $\C$. Thus $\alpha+\beta$, $\alpha\beta$, and $-\alpha$ are algebraic integers. Finally 0 and 1 satisfy $x$ and $x-1$. These facts verify all subring conditions.
